\documentclass[11pt]{article}
\usepackage[margin=1.1in]{geometry}
\usepackage{booktabs}
\usepackage{amsmath}
\usepackage{graphicx}
\usepackage{xcolor}
% Paper chrome uses the firm-microsim template blue (YGTemplate MyBlue);
% figures use the PolicyEngine palette separately.
\definecolor{paperblue}{rgb}{0,0,0.7}   % MyBlue, matching the sister paper
\definecolor{peteal}{HTML}{39C6C0}      % PolicyEngine teal accent
\usepackage{sectsty}
\allsectionsfont{\color{paperblue}}
\usepackage[colorlinks=true, citecolor=paperblue, linkcolor=paperblue, urlcolor=paperblue]{hyperref}
\usepackage[hang,flushmargin]{footmisc}
\usepackage{natbib}
\bibliographystyle{apalike}

% Robust figure inclusion: fall back to a placeholder if the PNG has not
% yet been generated by the analysis pipeline.
\newcommand{\safefig}[2][\linewidth]{%
  \IfFileExists{#2}{\includegraphics[width=#1]{#2}}{%
    \fbox{\parbox[c][0.35\linewidth][c]{0.9\linewidth}{\centering\color{gray}\small Figure pending: \texttt{\detokenize{#2}}}}}}

\title{\bfseries From council tax to land value tax:\\[0.15cm] who gains from a revenue-neutral reform?\thanks{\textbf{Acknowledgements:} We thank colleagues at PolicyEngine for valuable comments and helpful discussions.}\\[0.7cm]}
\author{Vahid Ahmadi\thanks{Research Associate at PolicyEngine. Email: \href{mailto:vahid@policyengine.org}{vahid@policyengine.org}.}
\and Max Ghenis\thanks{CEO, PolicyEngine. Email: \href{mailto:max@policyengine.org}{max@policyengine.org}.}}
\date{\small \ifcase\month\or January\or February\or March\or April\or May\or June\or July\or August\or September\or October\or November\or December\fi\ \the\year}

\begin{document}
\maketitle

\begin{abstract}
\noindent Council tax, the UK's principal recurrent property tax, is based on 1991 valuations and is regressive with respect to property value. We use PolicyEngine UK, an open-source static microsimulation model built on the Enhanced Family Resources Survey, to simulate replacing council tax's \pounds 57.6 billion of net revenue with a flat annual tax on the unimproved value of land. We impute land values to every household by combining Wealth and Assets Survey property wealth with region-specific land shares and an allocation of \pounds 2.06 trillion of corporate land, calibrating aggregates to an ONS National Balance Sheet total of \pounds 7.46 trillion for 2026--27. We calculate a budget-neutral rate of 0.77 per cent of land value per year. At this rate we estimate that the average household in each of income deciles 1--8 gains --- \pounds 481 per year (2.3 per cent of net income) in the poorest decile --- while deciles 9 and 10 lose \pounds 991 and \pounds 966 respectively, and 68 per cent of households gain overall. Absolute poverty before housing costs falls by 0.65 percentage points (1.04 points after housing costs). The Gini coefficient of equivalised household net income is essentially unchanged (+0.0004), yet when households are ranked by wealth the burden shifts markedly towards the wealthiest: the top wealth decile --- which we estimate holds 48.5 per cent of household land --- loses \pounds 5,573 per year on average and decile 9 loses \pounds 1,042, while wealth deciles 1--7 gain between \pounds 484 and \pounds 1,464. Our results suggest that a revenue-neutral land value tax would redistribute from land wealth towards households with little land while leaving conventional income-inequality statistics largely unaffected.
\end{abstract}

\medskip
\noindent\textbf{Keywords:} land value tax, council tax, microsimulation, tax incidence, poverty, inequality, United Kingdom

\newpage
\section{Introduction}
\label{sec:intro}

Council tax is the United Kingdom's principal recurrent tax on residential property, with gross receipts of approximately \pounds 62 billion forecast for 2026--27 \citep{obr2025}. Its design has attracted sustained criticism in the public-finance literature \citep{mirrlees2011, adam2020}. Dwellings in England and Scotland remain assessed on their estimated value as at 1 April 1991; the banding structure is capped, with the top band in England (Band H) covering all properties worth \pounds 320{,}000 or more \emph{in 1991 prices}, so that properties of widely differing current values share the top band and face nearly identical bills; and effective tax rates decline as a share of property value as that value rises. Council tax is therefore regressive with respect to the value of the asset it nominally taxes, bears a weak relationship to households' current ability to pay, and discourages investment in property improvements.

This paper examines the household-level consequences of abolishing council tax and replacing its net revenue, pound for pound, with a flat annual tax on the unimproved value of land. The land value tax (LVT) has a long intellectual pedigree, from \citet{george1879} through the optimal-taxation literature \citep{ramsey1927} to the review by \citet{mirrlees2011}, resting on the observation that the supply of land is fixed: a tax on pure land rents raises revenue without discouraging productive activity. The UK policy debate, however, has to date lacked a household-level quantification of the gains and losses such a reform would generate.

We provide such a quantification. Using PolicyEngine UK, an open-source static microsimulation model of the UK tax and benefit system built on an enhanced version of the Family Resources Survey \citep{policyengine}, we proceed in three steps. First, we impute a land value to every household by combining survey property wealth with region-specific land shares and an allocation of corporate-sector land, calibrating national aggregates to the ONS National Balance Sheet \citep{ons2025}. Second, we reconstruct council tax liabilities bottom-up by region, band and discount status, calibrated to Valuation Office Agency dwelling counts. Third, we simulate replacing the model's net council tax revenue of \pounds 57.6 billion (gross liabilities of \pounds 61.9 billion less \pounds 4.3 billion of council tax reduction) with a flat-rate LVT for fiscal year 2026--27; Section~\ref{sec:methodology} discusses how this model-based figure compares with the OBR receipts forecast. Against an estimated UK land base of \pounds 7.46 trillion, we calculate that revenue neutrality requires a rate of 0.77 per cent of land value per year.

Our results suggest three principal findings. First, the swap is progressive across most of the income distribution: we estimate that the average household in each of income deciles 1--8 gains, with the poorest decile gaining \pounds 481 per year --- 2.3 per cent of its net income --- while the top two deciles lose \pounds 991 and \pounds 966 per year respectively; 68 per cent of households gain overall. Second, the reform reduces poverty: the absolute poverty rate before housing costs falls by 0.65 percentage points, and by 1.04 points after housing costs. Third, the Gini coefficient of equivalised household net income is essentially unchanged (+0.0004) even though the burden shifts markedly towards the wealthy: ranked by household wealth rather than income, the top wealth decile loses \pounds 5,573 per year on average and decile 9 loses \pounds 1,042, while wealth deciles 1--7 gain between \pounds 484 and \pounds 1,464. These findings can be reconciled by observing that land wealth is far more unequally distributed than income --- we estimate that the top wealth decile holds 48.5 per cent of household land --- but is only weakly correlated with current income, so that a tax on land redistributes across the wealth distribution in ways that income-based inequality statistics register only faintly.

In addition to the central budget-neutral scenario, we trace out a rate grid from 0.5 to 5 per cent, documenting how revenue (from \pounds 37.3 billion to \pounds 373.2 billion), poverty and inequality vary with the rate, and we decompose the tax base into household and corporate land.

The remainder of the paper proceeds as follows. Section~\ref{sec:literature} situates the analysis in the literature on land taxation and UK property-tax reform. Section~\ref{sec:methodology} describes the data, the land-value imputation, the council tax baseline and the reform implementation. Section~\ref{sec:results} presents the distributional results. Section~\ref{sec:discussion} summarises the findings, discusses limitations and policy implications, and concludes. An appendix reports the full rate grid, a summary of landless households, a worked single-household example and a data and parameter summary.

\section{Related literature}
\label{sec:literature}

\subsection{The economics of land value taxation}

The case for taxing land rests on the classical observation that land is supplied inelastically. \citet{george1879}, in \emph{Progress and Poverty}, argued that the rental value of unimproved land is a socially created surplus --- driven by population, public infrastructure and agglomeration rather than by any effort of the owner --- and proposed a ``single tax'' on land values to replace taxes on labour and capital. The modern optimal-taxation literature formalises the intuition: because the base cannot shrink in response to the tax, a levy on pure land rents generates no substitution effect and hence no deadweight loss, making it the canonical example of efficient taxation in the \citet{ramsey1927} inverse-elasticity framework. \citet{feldstein1977} qualified the classical view by noting general-equilibrium portfolio effects, but the presumption that land taxes are unusually efficient survives.

Two further strands reinforce the theoretical case. \citet{arnott1979} establish the Henry George Theorem: under conditions of optimal city size, aggregate land rents exactly finance the efficient level of local public expenditure, giving land taxation a normative role beyond efficiency --- it is the natural fiscal counterpart of the public goods that generate land value. In cross-country empirical work, \citet{johansson2008} rank recurrent taxes on immovable property as the least harmful to long-run growth among the major tax instruments, ahead of consumption, personal income and corporate income taxes. More recently, \citet{bonnet2021} document that the post-1980 rise in wealth-to-income ratios is to a large extent a rise in the price of residential land, and argue that land is accordingly both an increasingly important and an administratively feasible tax base for contemporary economies.

The empirical evidence, while limited by the scarcity of implemented land taxes, is broadly consistent with the theory. \citet{oates1997} examine Pittsburgh's two-rate property tax --- which taxed land at more than five times the rate applied to structures --- and find that the city sustained substantially higher building activity than comparable cities in the same period; they attribute this partly to the land tax's neutrality with respect to development, while noting the difficulty of isolating the tax's contribution from concurrent local conditions. \citet{tideman1995} argues on theoretical grounds that land taxation is better than neutral, as it discourages speculative withholding of developable land. The contributions collected in \citet{dye2009} temper this optimism by documenting the practical assessment challenges --- above all, separating land from improvement value at scale --- that have historically limited adoption.

\subsection{UK property taxation and council tax}

The deficiencies of council tax are well documented. \citet{mirrlees2011} concluded that UK housing taxation lacks a coherent rationale: council tax is regressive in property value and based on obsolete 1991 valuations, while stamp duty land tax taxes transactions and thereby impedes mobility. The review proposed replacing both with a Housing Services Tax proportional to the consumption value of housing, complemented by a land value tax on business land --- an explicit endorsement of land taxation where valuation is feasible. Subsequent work has quantified the regressivity that motivates reform: \citet{adam2020} show that effective council tax rates as a share of property value fall sharply across the price distribution, and that the absence of revaluation in England and Scotland leaves relative price changes since 1991 --- most notably London's appreciation --- untaxed. \citet{muellbauer2005} argues that reforming council tax into a proportional property tax with regular automatic revaluation would also stabilise the housing market, and the OECD's most recent survey of the UK likewise recommends updating property valuations and reforming council tax \citep{oecd2024}.

A body of applied UK work has modelled alternatives. \citet{corlett2018} examine options including a proportional property tax and a partial LVT at the aggregate and regional level; \citet{fairershare2021} propose a proportional property tax of 0.48 per cent of current property value and estimate that roughly three-quarters of households would pay less --- a useful comparator for our finding that 68 per cent of households gain at a 0.77 per cent tax on the narrower land base. \citet{ippr2019} model council tax reform for London. In Scotland, the cross-party Commission on Local Tax Reform \citep{scotgov2015} concluded that the present system must end and assessed property, land and income bases, and \citet{hughes2018} examine LVT options for the Scottish Land Commission, highlighting valuation and transition as the binding constraints. \citet{leunig2024} propose replacing both council tax and stamp duty with proportional levies on current values. Closest to our exercise, \citet{apgwilym2020} assess the feasibility of a local land value tax for Wales for the Welsh Government, simulating revenue-neutral rates at local-authority level using small-area land value estimates.

\subsection{Contribution}

Relative to this literature, our contribution is a household-level microsimulation of a full council-tax-to-LVT swap for the UK, with land values imputed to every household in a nationally representative micro-dataset, benchmarked against national balance-sheet aggregates, and evaluated within a complete model of the tax and benefit system. Prior UK LVT analyses have generally worked with regional, band-level or local-authority aggregates --- the Welsh assessment of \citet{apgwilym2020} being the most detailed to date --- rather than with household microdata. The microdata buy us three things that aggregate analyses cannot deliver: the joint distribution of land and income, which drives the winners-and-losers results and the divergence between income-ranked and wealth-ranked incidence; poverty rates measured against the statutory equivalised, housing-cost-adjusted thresholds; and a decomposition of who bears the corporate component of the base. We note in Section~\ref{sec:results} that the swap turns out to generate no interactions with means-tested benefit entitlements, so the household-level net income change is exactly the difference between the abolished council tax bill and the new land tax bill --- a finding, not an assumption, and one that materially simplifies interpretation. To our knowledge this is the first open-source, household-level distributional analysis of a revenue-neutral LVT for the whole of the UK, and the code and data pipeline are fully reproducible \citep{policyengine}.

\section{Data and methodology}
\label{sec:methodology}

\subsection{Microdata}

Our simulations use PolicyEngine UK \citep{policyengine}, an open-source static microsimulation model of the UK tax and benefit system, run on the Enhanced Family Resources Survey (EFRS) for 2023--24, uprated to fiscal year 2026--27. The EFRS augments the standard Family Resources Survey in two ways. First, because household surveys under-represent the top of the income distribution, synthetic high-income records are generated by training a quantile random forest on HMRC's Survey of Personal Incomes (SPI) and imputing the upper tail of incomes. Second, household weights are calibrated by gradient descent so that the weighted dataset matches a large set of national and regional administrative totals from the ONS, HMRC and DWP --- including population by age and region, income components, and benefit caseloads. Wealth variables, including property wealth and corporate (financial) wealth, are imputed from the ONS Wealth and Assets Survey (WAS).

\subsection{Council tax baseline}

The FRS records council tax imperfectly, so we reconstruct liabilities. Each household is assigned a council tax band (A--H in England and Scotland, A--I in Wales) and a bill equal to the average band rate in its region, with the 25 per cent single-person discount imputed to single-adult households. Band assignment probabilities are calibrated so that the distribution of dwellings by region and band matches Valuation Office Agency (and Scottish Assessors) counts. This bottom-up reconstruction yields, for 2026--27, gross council tax of \pounds 61.9 billion; netting off \pounds 4.3 billion of means-tested council tax reduction gives net revenue of
\begin{equation}
R_{\mathrm{CT}} \;=\; 61.9 - 4.3 \;=\; \pounds 57.6 \text{ billion},
\label{eq:ctrev}
\end{equation}
which is the amount the land value tax must replace. We note a calibration caveat: this model-based figure exceeds the Office for Budget Responsibility's accruals-based council tax receipts forecast of approximately \pounds 53.7 billion \citep{obr2025}. Were the replacement target instead calibrated to OBR-consistent receipts, the revenue-neutral rate would be approximately 0.71 rather than 0.77 per cent; all results below use the internally consistent model figure, which ensures that the simulated LVT replaces exactly the council tax the model abolishes household by household. Northern Ireland operates a separate domestic rates system, which we leave unchanged (see Section~\ref{sec:discussion}).

\subsection{Imputing land values}
\label{sec:landvalues}

The tax base is the unimproved value of land held directly (residential land under owner-occupied and other household-owned property) and indirectly (land on corporate balance sheets, attributed to households as ultimate owners). For household $i$ in region $r$, land value is
\begin{equation}
L_i \;=\; \underbrace{\lambda_r \, P_i}_{\text{household land}} \;+\; \underbrace{\frac{K_i}{\sum_j K_j w_j} \, L^{\mathrm{corp}}}_{\text{allocated corporate land}},
\label{eq:land}
\end{equation}
where $P_i$ is WAS-imputed property wealth, $\lambda_r$ is the region-specific land share of property value, $K_i$ is the household's corporate wealth, $w_j$ are survey weights, and $L^{\mathrm{corp}} = \pounds 2.06$ trillion is aggregate corporate-sector land. Land shares $\lambda_r$ vary from 85 per cent in London to 42 per cent in the North East, reflecting the fact that in high-price regions almost all of the price differential is location value rather than structure cost. Corporate land is allocated proportionally to each household's corporate wealth, treating shareholdings (direct and via pensions) as the ultimate claim on corporate assets.

Aggregates are calibrated to targets in the PolicyEngine UK data pipeline derived from the sectoral tables underlying the ONS National Balance Sheet \citep{ons2025}: UK land of \pounds 7.10 trillion at end-2024, split into \pounds 5.04 trillion household and \pounds 2.06 trillion corporate. (These derived targets differ from the headline bulletin presentation, which reports household non-produced assets of approximately \pounds 4.6 trillion in 2024 on a different sectoral grouping.) We uprate to 2026--27 using OBR per-capita nominal GDP growth \citep{obr2025}, giving a total land base of
\begin{equation}
L \;=\; \sum_i L_i \, w_i \;=\; \pounds 7.46 \text{ trillion},
\end{equation}
comprising \pounds 5.40 trillion of household land and \pounds 2.06 trillion of corporate land. The model total is 105.1 per cent of the (un-uprated) 2024 ONS benchmark. Land so measured is 68.4 per cent of total property wealth and 27.7 per cent of total household wealth (\pounds 26.98 trillion); mean land value per household is \pounds 232{,}672 and the median is \pounds 93{,}796.

\subsection{The reform and the revenue-neutral rate}

The reform (i) abolishes council tax in Great Britain and (ii) levies a flat annual tax at rate $r$ on each household's land value $L_i$, implemented in PolicyEngine UK via the parameters \texttt{gov.contrib.ubi\_center.land\_value\_tax.rate} and \texttt{gov.contrib.abolish\_council\_tax}. The budget-neutral rate equates LVT revenue with net council tax revenue:
\begin{equation}
r^{*} \;=\; \frac{R_{\mathrm{CT}}}{L} \;=\; \frac{57.6}{7{,}463} \;=\; 0.77\%.
\label{eq:rate}
\end{equation}
The net income change for household $i$ is $\Delta y_i = \mathrm{CT}_i - r^{*} L_i$, where $\mathrm{CT}_i$ is the abolished council tax liability (net of council tax reduction). Three maintained assumptions govern incidence, all standard in static microsimulation. First, incidence is assigned entirely to current landowners: we assume no behavioural response and no capitalisation of the tax into land prices. Second, rents are held fixed, so renters --- who own no residential land --- retain their abolished council tax liability in full and pay LVT only on any corporate wealth they hold. Third, the land base is held fixed across the rate grid, consistent with the inelastic supply of land. All results are for simulation year 2026, with poverty measured as absolute poverty against a baseline-fixed threshold, before housing costs (BHC) and after housing costs (AHC), and inequality by the Gini coefficient of equivalised household net income. In addition to the central rate $r^{*}=0.77\%$ we simulate a grid of rates $\{0.5, 1.0, 1.5, 2.0, 3.0, 5.0\}$ per cent to trace out revenue and distributional effects, and decompose revenue by base scope (household-only versus corporate-only land).

\section{Results}
\label{sec:results}

\subsection{Where land value sits}
\label{sec:whereland}

The distributional consequences of an LVT are determined by the distribution of its base, which we describe first. Table~\ref{tab:land_income} reports average land value, property wealth and land shares by decile of equivalised household net income. Two features stand out. First, average land value rises with income, but weakly and unevenly: it is \pounds 95{,}302 in the poorest decile and \pounds 442{,}935 in the richest, a ratio of 4.6 to 1 against an income ratio of 6.5 to 1, and the profile is far from monotonic --- decile 2 (\pounds 217{,}644) holds nearly twice as much land on average as decile 3 (\pounds 115{,}702). Some of this reflects genuine heterogeneity, principally asset-rich retired households with low current income; some reflects imputation noise in a base constructed from survey property wealth. Either way, the income gradient of the tax base is shallow, which is what drives the results in Section~\ref{sec:poverty}. Second, the top two income deciles together hold 35.5 per cent of land, whereas, as Section~\ref{sec:wealth_impact} shows, the top two \emph{wealth} deciles hold 66.6 per cent.

Table~\ref{tab:land_income} separates the two components of the base, which are not comparable quantities: household land is a share of the household's own property wealth, while allocated corporate land is an attributed claim on corporate balance sheets that households holding no property may still bear. Reporting them together against property wealth --- as a single ``land value'' column would --- can make land appear to exceed the property wealth it is derived from.

\begin{table}[htbp]
\centering
\caption{Land value and property wealth by decile of equivalised household net income, 2026--27}
\label{tab:land_income}
\begin{tabular}{lrrrrrr}
\toprule
Income & Avg.\ income & Avg.\ household & Avg.\ corporate & Avg.\ total & Avg.\ property & Share of \\
decile & (\pounds/yr) & land (\pounds) & land (\pounds) & land (\pounds) & wealth (\pounds) & land (\%) \\
\midrule
1  & 21{,}701  & 75{,}848  & 19{,}454  & 95{,}302  & 106{,}712 & 4.7 \\
2  & 31{,}068  & 171{,}344 & 46{,}300  & 217{,}644 & 198{,}309 & 8.8 \\
3  & 37{,}182  & 89{,}340  & 26{,}362  & 115{,}702 & 144{,}651 & 5.0 \\
4  & 38{,}400  & 152{,}099 & 50{,}164  & 202{,}263 & 204{,}378 & 9.3 \\
5  & 43{,}127  & 144{,}852 & 52{,}305  & 197{,}157 & 213{,}009 & 9.0 \\
6  & 53{,}985  & 160{,}086 & 50{,}458  & 210{,}544 & 231{,}152 & 8.5 \\
7  & 55{,}357  & 138{,}281 & 48{,}129  & 186{,}410 & 221{,}029 & 8.1 \\
8  & 69{,}125  & 191{,}548 & 88{,}673  & 280{,}220 & 308{,}960 & 10.9 \\
9  & 77{,}096  & 294{,}523 & 127{,}118 & 421{,}641 & 451{,}521 & 17.5 \\
10 & 141{,}327 & 294{,}942 & 147{,}992 & 442{,}935 & 429{,}497 & 18.0 \\
\bottomrule
\end{tabular}

\medskip
{\footnotesize\emph{Source:} Authors' calculations using PolicyEngine UK.}
\end{table}

Regional variation is at least as pronounced as variation by income. Average household land value is \pounds 369{,}802 in London and \pounds 286{,}311 in the South East, compared with \pounds 153{,}631 in the North East and \pounds 159{,}535 in Scotland (Figure~\ref{fig:land_region}) --- a ratio of 2.4 to 1 that reflects both higher property prices and higher land shares of those prices in the South. Across household types, pensioner couples hold the most land on average (\pounds 431{,}377), followed by working-age couples without children (\pounds 303{,}440); lone parents, who are predominantly renters, hold very little (\pounds 7{,}797) (Figure~\ref{fig:land_famtype}). Ranked by total wealth rather than income, concentration is far greater: the bottom three wealth deciles hold essentially no land, while the top wealth decile averages \pounds 1{,}088{,}892 --- 47.4 per cent of all land (Table~\ref{tab:wealth_impact} below). Overall, 17.2 per cent of households hold no land.

\begin{figure}[htbp]
\centering
\safefig{../results/figures/fig2_land_by_income_decile.png}
\caption{Average land value by income decile.}
\label{fig:land_income}
\end{figure}

\begin{figure}[htbp]
\centering
\safefig{../results/figures/fig3_land_by_region.png}
\caption{Average household land value by region.}
\label{fig:land_region}
\end{figure}

\begin{figure}[htbp]
\centering
\safefig{../results/figures/fig4_land_by_family_type.png}
\caption{Average household land value by family type.}
\label{fig:land_famtype}
\end{figure}

\subsection{Revenue and the rate schedule}
\label{sec:revenue}

Because the base is fixed in the static framework, LVT revenue is linear in the rate: each 0.1 percentage point raises approximately \pounds 7.5 billion. Table~\ref{tab:revenue} reports revenue across the rate grid. A 0.5 per cent rate raises \pounds 37.3 billion --- \pounds 20.3 billion less than council tax --- while a 1 per cent rate raises \pounds 74.6 billion and a 5 per cent rate raises \pounds 373.2 billion. The budget-neutral rate is 0.77 per cent (equation~\ref{eq:rate}). Decomposing the base at a 1 per cent rate, household land contributes \pounds 54.0 billion and corporate land \pounds 20.6 billion of the \pounds 74.6 billion total. The corporate component --- 28 per cent of revenue --- falls on households in proportion to their shareholdings, and its inclusion permits a lower rate for any given revenue target than a purely residential base would require.

\begin{table}[htbp]
\centering
\caption{LVT revenue by rate, 2026--27}
\label{tab:revenue}
\begin{tabular}{lrrrr}
\toprule
Rate & LVT revenue & Council tax & Net revenue & Avg.\ LVT per \\
     & (\pounds bn) & (\pounds bn) & (\pounds bn) & household (\pounds) \\
\midrule
0.5\% & 37.3  & 57.6 & $-$20.3 & 1{,}163 \\
1.0\% & 74.6  & 57.6 & 17.0    & 2{,}327 \\
1.5\% & 111.9 & 57.6 & 54.3    & 3{,}490 \\
2.0\% & 149.3 & 57.6 & 91.6    & 4{,}653 \\
3.0\% & 223.9 & 57.6 & 166.3   & 6{,}980 \\
5.0\% & 373.2 & 57.6 & 315.5   & 11{,}634 \\
\bottomrule
\end{tabular}

\medskip
{\footnotesize\emph{Note:} Net revenue is LVT revenue less replaced council tax revenue. \emph{Source:} Authors' calculations using PolicyEngine UK.}
\end{table}

\begin{figure}[htbp]
\centering
\safefig{../results/figures/fig7_revenue_by_rate.png}
\caption{LVT revenue by rate against the council tax replacement target.}
\label{fig:revenue}
\end{figure}

\subsection{Distributional impact by income decile}
\label{sec:income_impact}

Table~\ref{tab:income_impact} and Figure~\ref{fig:net_income} report the average change in household net income at the budget-neutral 0.77 per cent rate, by income decile. Because the swap generates no benefit interactions --- a result we verify in Section~\ref{sec:robustness} --- the change is exactly the abolished council tax bill less the land tax bill, $\Delta y_i = \mathrm{CT}_i - r^{*}L_i$.

The pattern across income deciles is best described as flat with a step down at the top. The poorest decile gains most in both absolute and relative terms, \pounds 647 per year or 3.0 per cent of net income, because its average council tax bill of \pounds 1{,}383 --- inflated relative to its means by the compressed banding structure --- is nearly twice the \pounds 736 land tax charged on its modest holdings. Deciles 3 and 7 gain \pounds 578 and \pounds 455; deciles 5, 6 and 8 are close to unchanged; and deciles 2 and 4 record small average \emph{losses} of \pounds 163 and \pounds 60, reflecting the uneven land profile documented in Table~\ref{tab:land_income} rather than anything systematic about those income levels. Only in deciles 9 and 10 is the loss substantial and clearly attributable to income: \pounds 1{,}016 and \pounds 963 per year, as land holdings of \pounds 420{,}000--\pounds 443{,}000 attract liabilities of \pounds 3{,}250--\pounds 3{,}420 against council tax savings of \pounds 2{,}240--\pounds 2{,}460. In percentage terms the loss is larger for decile 9 ($-1.3$ per cent) than for decile 10 ($-0.7$ per cent), because top-decile incomes are considerably higher while average land holdings are only slightly larger.

We stress that the middle of this profile is not precisely estimated. The differences between adjacent deciles in the range 2--8 are of the order of a few hundred pounds a year on average net incomes of \pounds 31{,}000--\pounds 69{,}000, and they rest on imputed land values whose decile means are visibly noisy. The result that should be read from Table~\ref{tab:income_impact} is that the reform is roughly distributionally neutral across the bottom eight deciles taken together and clearly costly to the top two --- not that any particular middle decile gains or loses.

\begin{table}[htbp]
\centering
\caption{Impact of the 0.77\% budget-neutral LVT by income decile}
\label{tab:income_impact}
\begin{tabular}{lrrrrrr}
\toprule
Income & Avg.\ LVT & Avg.\ CT & Avg.\ net & \% of net & Winners & Losers \\
decile & (\pounds) & saved (\pounds) & change (\pounds) & income & (\%) & (\%) \\
\midrule
1  & 736     & 1{,}383 & $+$647    & $+$3.0 & 79.0 & 18.3 \\
2  & 1{,}681 & 1{,}518 & $-$163    & $-$0.5 & 72.0 & 26.9 \\
3  & 894     & 1{,}472 & $+$578    & $+$1.6 & 76.5 & 22.4 \\
4  & 1{,}562 & 1{,}502 & $-$60     & $-$0.2 & 72.6 & 26.6 \\
5  & 1{,}523 & 1{,}668 & $+$145    & $+$0.3 & 69.1 & 29.7 \\
6  & 1{,}626 & 1{,}805 & $+$179    & $+$0.3 & 70.9 & 28.6 \\
7  & 1{,}440 & 1{,}895 & $+$455    & $+$0.8 & 76.0 & 23.7 \\
8  & 2{,}164 & 2{,}199 & $+$35     & $+$0.1 & 67.9 & 30.9 \\
9  & 3{,}256 & 2{,}240 & $-$1{,}016 & $-$1.3 & 48.1 & 51.4 \\
10 & 3{,}421 & 2{,}458 & $-$963    & $-$0.7 & 48.6 & 50.8 \\
\bottomrule
\end{tabular}

\medskip
{\footnotesize\emph{Note:} ``CT saved'' denotes the abolished council tax liability net of council tax reduction. Winners and losers do not sum to 100 per cent because a small share of households are unaffected. \emph{Source:} Authors' calculations using PolicyEngine UK.}
\end{table}

\begin{figure}[htbp]
\centering
\safefig{../results/figures/fig1_net_change_by_income_decile.png}
\caption{Average net income change by income decile at the 0.77\% budget-neutral rate.}
\label{fig:net_income}
\end{figure}

\subsection{Distributional impact by wealth decile}
\label{sec:wealth_impact}

Table~\ref{tab:wealth_impact} and Figure~\ref{fig:net_wealth} report the corresponding results when households are ranked by total wealth rather than income.\footnote{We construct equal-weight wealth deciles directly. PolicyEngine UK's own \texttt{household\_wealth\_decile} variable is degenerate at the bottom of the distribution --- the mass of households with zero or negative net wealth is tied, leaving decile one empty and placing 22.6 per cent of households in decile two --- which would make the decile averages in this table incomparable.} The gradient here is monotonic and steep, in sharp contrast to the income ranking. The bottom two wealth deciles own essentially no land and retain their council tax savings in full, gaining \pounds 898 and \pounds 902 per year; deciles 3--7 gain between \pounds 523 and \pounds 1{,}480; decile 8 loses \pounds 268; decile 9 loses \pounds 1{,}069; and the top wealth decile --- averaging \pounds 1.09 million of land --- loses \pounds 5{,}752 per year, or 8.8 per cent of its net income, with an average LVT liability of \pounds 8{,}410 against council tax savings of \pounds 2{,}658. In the top wealth decile 98.4 per cent of households lose; in each of the bottom seven, more than 82 per cent gain. The incidence of the reform tracks land wealth almost mechanically, which is exactly what the instrument is designed to do.\footnote{Revenue neutrality holds in aggregate --- LVT revenue equals abolished council tax revenue at \pounds 57.6 billion --- so the population-weighted average net-income change is zero. It is identical whether households are grouped by income or by wealth decile, since the two tables show the same households re-sorted; the per-decile averages need not sum to zero because decile group sizes are equal in population but not in the underlying quantities.}

\begin{table}[htbp]
\centering
\caption{Impact of the 0.77\% budget-neutral LVT by wealth decile}
\label{tab:wealth_impact}
\begin{tabular}{lrrrrrr}
\toprule
Wealth & Avg.\ land & Avg.\ LVT & Avg.\ CT & Avg.\ net & Winners & Losers \\
decile & value (\pounds) & (\pounds) & saved (\pounds) & change (\pounds) & (\%) & (\%) \\
\midrule
1  & 1             & 0       & 898     & $+$898     & 82.1 & 13.7 \\
2  & 29            & 0       & 902     & $+$902     & 83.2 & 11.4 \\
3  & 1{,}265       & 10      & 1{,}177 & $+$1{,}168 & 91.2 & 8.1 \\
4  & 25{,}263      & 195     & 1{,}675 & $+$1{,}480 & 94.4 & 5.6 \\
5  & 76{,}699      & 592     & 1{,}833 & $+$1{,}241 & 91.8 & 8.2 \\
6  & 138{,}998     & 1{,}074 & 2{,}027 & $+$953     & 89.5 & 10.5 \\
7  & 212{,}347     & 1{,}640 & 2{,}163 & $+$523     & 82.9 & 17.1 \\
8  & 321{,}469     & 2{,}483 & 2{,}215 & $-$268     & 48.0 & 52.0 \\
9  & 451{,}386     & 3{,}486 & 2{,}417 & $-$1{,}069 & 20.0 & 80.0 \\
10 & 1{,}088{,}892 & 8{,}410 & 2{,}658 & $-$5{,}752 & 1.6  & 98.4 \\
\bottomrule
\end{tabular}

\medskip
{\footnotesize\emph{Note:} Deciles of total household wealth. \emph{Source:} Authors' calculations using PolicyEngine UK.}
\end{table}

\begin{figure}[htbp]
\centering
\safefig{../results/figures/fig5_net_change_by_wealth_decile.png}
\caption{Average net income change by wealth decile at the 0.77\% budget-neutral rate.}
\label{fig:net_wealth}
\end{figure}

\subsection{Winners and losers within deciles}
\label{sec:winners}

Decile averages conceal substantial within-decile variation, because land holdings vary widely within any given income band: an outright owner-occupier in London and a private renter in the North East may occupy the same income decile. At the budget-neutral rate, 68.2 per cent of all households gain. The share of gaining households declines only gradually with income through deciles 1--8 (from 79.0 to 67.9 per cent) before falling to just under one half in deciles 9--10 (Figure~\ref{fig:winners}). Even in the top income decile, 48.6 per cent of households gain --- predominantly high earners who rent or own property of modest value --- while 18.3 per cent of bottom-decile households lose, typically low-income owners of valuable land. Sorted by wealth, the gradient is considerably steeper: the share of gaining households exceeds 82 per cent in each of the bottom seven wealth deciles and falls to 1.6 per cent in the top decile. The 17.2 per cent of households holding no land gain \pounds 832 per year on average, and 81.5 per cent of them gain; the remainder are largely households whose council tax was already fully offset by council tax reduction.

That the winner share is high while the average income-decile effects are small is not a contradiction. The reform moves a fixed \pounds 57.6 billion from a base that is broad and shallow (council tax, capped at the top) to one that is narrow and deep (land, highly concentrated). Most households pay a little less; a small minority pay a great deal more. This is visible in the wealth ranking and almost invisible in the income ranking.

\begin{figure}[htbp]
\centering
\safefig{../results/figures/fig6_winners_losers_by_decile.png}
\caption{Share of winners and losers by income decile at the 0.77\% budget-neutral rate.}
\label{fig:winners}
\end{figure}

\subsection{Poverty and inequality}
\label{sec:poverty}

Table~\ref{tab:poverty} reports poverty and inequality across the rate grid, measured against baseline-fixed absolute poverty thresholds. At the budget-neutral 0.77 per cent rate the two poverty measures move in opposite directions. Absolute poverty after housing costs falls from 14.75 to 14.04 per cent, a reduction of 0.70 percentage points, while poverty before housing costs is essentially unchanged, edging up from 11.53 to 11.58 per cent ($+0.04$ points). The Gini coefficient of equivalised household net income is likewise unchanged, rising by 0.0003 from a baseline of 0.3641.

The BHC and AHC divergence is informative rather than a puzzle. The AHC measure nets housing costs out of income, and the households lifted above the AHC line are disproportionately low-income renters, who gain their entire council tax bill and owe almost no land tax. The BHC measure does not, so it gives more weight to low-income owner-occupiers, who lose. The near-zero BHC effect is the sum of these two flows, and it is small enough that we would not claim a poverty reduction on the BHC measure at all: the honest statement is that the revenue-neutral swap reduces after-housing-cost poverty modestly and leaves before-housing-cost poverty unchanged.

The grid shows how sensitive this is to the rate, and therefore to the replacement target. At 0.5 per cent --- which is not revenue-neutral, reducing net revenue by \pounds 20.3 billion --- poverty falls on both measures (BHC $-0.60$, AHC $-0.91$ points) and the Gini declines slightly ($-0.0028$). At 1 per cent, which raises \pounds 17.0 billion of net revenue, BHC poverty rises 0.83 points while AHC poverty is still 0.30 points below baseline. At higher rates poverty and measured income inequality rise substantially --- at 5 per cent, BHC poverty rises by 12.0 points and the Gini by 0.10 --- because the static analysis levies cash charges on asset-rich, income-poor households with no offsetting use of the additional revenue. These upper-grid rows are the gross incidence of the tax alone and not complete reform packages: any practical proposal at those rates would recycle the net revenue, and the rows should not be read as forecasts of what such a package would do. The wealth Gini (0.7038) is mechanically unchanged, since the simulation does not model asset-price responses; Section~\ref{sec:robustness} reports what capitalisation would imply for it.

\begin{table}[htbp]
\centering
\caption{Poverty and inequality by LVT rate (council tax abolished throughout)}
\label{tab:poverty}
\begin{tabular}{lrrrrr}
\toprule
Rate & Poverty BHC & $\Delta$ BHC & Poverty AHC & $\Delta$ AHC & $\Delta$ Gini \\
     & (\%) & (pp) & (\%) & (pp) & \\
\midrule
Baseline & 11.53 & ---      & 14.75 & ---      & --- \\
0.5\%    & 10.94 & $-$0.60  & 13.84 & $-$0.91  & $-$0.0028 \\
0.77\%   & 11.58 & $+$0.04  & 14.04 & $-$0.70  & $+$0.0003 \\
1.0\%    & 12.36 & $+$0.83  & 14.45 & $-$0.30  & $+$0.0033 \\
1.5\%    & 13.82 & $+$2.29  & 15.62 & $+$0.87  & $+$0.0108 \\
2.0\%    & 15.55 & $+$4.02  & 17.56 & $+$2.81  & $+$0.0197 \\
3.0\%    & 18.86 & $+$7.32  & 21.52 & $+$6.78  & $+$0.0417 \\
5.0\%    & 23.70 & $+$12.17 & 26.61 & $+$11.86 & $+$0.1013 \\
\bottomrule
\end{tabular}

\medskip
{\footnotesize\emph{Note:} Absolute poverty against baseline-fixed thresholds, before (BHC) and after (AHC) housing costs. Gini refers to equivalised household net income; baseline value 0.3641. Rates above 0.77\% raise net revenue that is not recycled in the simulation. \emph{Source:} Authors' calculations using PolicyEngine UK.}
\end{table}

\begin{figure}[htbp]
\centering
\safefig{../results/figures/fig8_poverty_gini_by_rate.png}
\caption{Poverty and Gini changes across the LVT rate grid.}
\label{fig:povgini}
\end{figure}

\subsection{Validation and robustness}
\label{sec:robustness}

\paragraph{The swap has no benefit interactions.} Running the reform through the full tax and benefit model and comparing with the arithmetic $\Delta y_i = \mathrm{CT}_i - r L_i$, we find the two agree to within \pounds 0.25 for every household at every rate tested (0.5, 1.0 and 2.0 per cent), and the share of households for which they differ by more than \pounds 1 is zero. Abolishing council tax and levying a land tax changes no means-tested entitlement in the UK system: council tax reduction is itself abolished with the tax, and neither instrument enters the income or capital tests for Universal Credit, Pension Credit or the legacy benefits. The consequence is that the microdata, not the tax-benefit calculator, do all the work in this reform --- the calculator's contribution is the council tax baseline, the equivalisation and the statutory poverty thresholds. We regard this as worth stating explicitly because it is easy to assume otherwise, and it means the results can be reproduced and audited without re-running the model.

\paragraph{Robustness.} Table~\ref{tab:robustness} varies the replacement target, the composition of the base and the instrument. The qualitative results are stable. The share of households gaining stays within a 65--70 per cent band across every specification. The budget-neutral rate ranges from 0.72 to 1.07 per cent, the top of that range being the household-land-only base, which is 28 per cent smaller. Calibrating the replacement target to OBR receipts rather than modelled liabilities lowers the rate to 0.72 per cent and, mechanically, makes every group better off --- it is a \pounds 3.9 billion net tax cut rather than a neutral swap, and we therefore treat it as a bound rather than a competing central case. Excluding Northern Ireland, whose households would otherwise pay the land tax on top of domestic rates, raises the rate to 0.79 per cent and leaves the distribution essentially unchanged. Assuming half of corporate equity is foreign-owned and therefore outside the household base raises the required rate to 0.90 per cent, the largest single sensitivity in the table, but still leaves 68.0 per cent of households gaining. Perturbing every regional land share by $\pm 10$ per cent, or rescaling household land to the un-uprated ONS 2024 benchmark, changes the rate by at most 0.06 percentage points and the decile effects by tens of pounds.

The final row is a comparator rather than a robustness check: a proportional tax on total property value, structures included, at the 0.73 per cent rate that raises the same revenue. It is less progressive on the wealth ranking --- the top wealth decile loses \pounds 4{,}217 rather than \pounds 5{,}752 --- and it produces fewer winners (64.4 against 68.2 per cent), because the base is spread more evenly across the property-owning population instead of being concentrated where location value is highest. That difference is the quantitative content of the land/structures distinction.

\begin{table}[htbp]
\centering
\footnotesize
\caption{Robustness: revenue-neutral rate and distributional summary under alternative assumptions}
\label{tab:robustness}
\begin{tabular}{lrrrrrrr}
\toprule
 & Base & Rate & Winners & Decile 1 & Decile 10 & Wealth 10 & $\Delta$ Poverty \\
Specification & (\pounds tn) & (\%) & (\%) & (\pounds) & (\pounds) & (\pounds) & AHC (pp) \\
\midrule
Central                              & 7.46 & 0.77 & 68.2 & $+$647 & $-$963    & $-$5{,}752 & $-$0.70 \\
OBR receipts target                  & 7.46 & 0.72 & 69.1 & $+$698 & $-$730    & $-$5{,}177 & $-$0.73 \\
Great Britain only                   & 7.34 & 0.79 & 67.9 & $+$639 & $-$1{,}016 & $-$5{,}747 & $-$0.56 \\
Household land only                  & 5.40 & 1.07 & 69.6 & $+$574 & $-$688    & $-$5{,}364 & $-$0.56 \\
Half of corporate land foreign-owned & 6.43 & 0.90 & 68.0 & $+$617 & $-$848    & $-$5{,}589 & $-$0.59 \\
Household land scaled to ONS 2024    & 7.10 & 0.81 & 68.2 & $+$651 & $-$977    & $-$5{,}772 & $-$0.71 \\
Land shares $-10$\%                  & 6.92 & 0.83 & 68.2 & $+$653 & $-$985    & $-$5{,}782 & $-$0.70 \\
Land shares $+10$\%                  & 8.00 & 0.72 & 68.2 & $+$642 & $-$945    & $-$5{,}726 & $-$0.56 \\
\midrule
\emph{Comparator:} proportional property tax & 7.91 & 0.73 & 64.4 & $+$605 & $-$674 & $-$4{,}217 & $-$0.76 \\
\bottomrule
\end{tabular}

\medskip
{\footnotesize\emph{Note:} Each row sets the rate so that revenue equals the replacement target on that row's base. ``Decile 1'' and ``Decile 10'' are average net income changes by income decile; ``Wealth 10'' is the top wealth decile. \emph{Source:} Authors' calculations using PolicyEngine UK; \texttt{analysis/extensions.py}.}
\end{table}

\paragraph{Capitalisation.} The central results are flow accounting and assume no capitalisation. If instead the levy is fully anticipated and permanent, it capitalises into land prices: a rate $r$ discounted at $\rho$ implies a proportional fall of $r/\rho$ in land values. At $\rho = 5$ per cent the 0.77 per cent rate implies a 15.4 per cent fall in land prices and an aggregate wealth loss of \pounds 1.15 trillion; at $\rho = 3$ per cent, a 25.7 per cent fall and \pounds 1.92 trillion. This is a transfer to the state of an asset value that was itself partly the capitalisation of past public investment, and under the textbook assumption it lands overwhelmingly on existing owners at announcement rather than on subsequent purchasers. Whether it does so in practice is contested: \citet{hoj2018} find full capitalisation in Danish data, whereas \citet{nielsson2024} cannot reject zero and argue the burden is shared with tenants and future buyers. We report the arithmetic under full capitalisation because it bounds the wealth effect, not because the evidence compels it. Table~\ref{tab:capitalisation} in the appendix reports the implied one-off wealth loss by wealth decile: at $\rho=3$ per cent the top wealth decile loses \pounds 280{,}331 of land value, 6.5 per cent of its total wealth, against \pounds 323 for decile 3. Two points follow. First, the wealth-decile gradient of the reform is steeper still on a stock than on a flow basis. Second, the flow results in Sections~\ref{sec:income_impact}--\ref{sec:wealth_impact} should be read as impact incidence in the first years, before the announcement effect on prices works through.

\section{Discussion}
\label{sec:discussion}

\subsection{Summary of findings}

Our simulations indicate that a revenue-neutral replacement of council tax with a 0.77 per cent flat land value tax would leave 68 per cent of households better off, reduce absolute poverty after housing costs by 0.70 percentage points, and shift the tax burden substantially up the wealth distribution, while leaving before-housing-cost poverty and the Gini coefficient of equivalised household net income essentially unchanged. The dissociation between the wealth-decile and income-Gini results is the paper's main analytical point, and it generalises: distributional statistics computed over the income distribution can be nearly uninformative about a reform whose incidence is organised by wealth. An LVT is a tax on a stock that is highly concentrated --- we estimate that the top wealth decile holds 47.4 per cent of all land, and the wealth Gini is 0.70 against 0.36 for income --- assessed on households whose position in the income distribution is only weakly related to that stock. A reader who evaluated this reform only on the Gini and the BHC poverty rate, the two statistics that dominate UK distributional debate, would conclude that it does almost nothing. Table~\ref{tab:wealth_impact} shows that it does a great deal.

We are correspondingly cautious about the income-decile profile. Once the top two deciles are set aside, the average effects in deciles 2--8 are a few hundred pounds a year on incomes of \pounds 31{,}000--\pounds 69{,}000, they are not monotonic in income, and they rest on imputed land values whose decile means are visibly noisy. We do not think the data support claims about which particular middle decile gains; they support the claim that the reform is roughly neutral across the bottom eight deciles and costly to the top two.

The efficiency case, which our static framework does not quantify, points in the same direction as the equity case. Whereas council tax is levied on the value of dwellings, including improvements, an LVT imposes no charge on constructing, extending or upgrading property, and its base is immobile \citep{george1879, oates1997}.

\subsection{Limitations}

Several caveats qualify the results.

\paragraph{Static incidence.} We assign the full burden to current landowners, with no behavioural response and no capitalisation in the central results. Standard theory predicts that an anticipated recurrent land tax capitalises into lower land prices \citep{feldstein1977}, so part of the burden would fall on owners as a one-off wealth loss at announcement rather than as a permanent flow; conversely, abolishing council tax would be expected to capitalise into higher property prices. The empirical literature does not settle the question: \citet{hoj2018} find full capitalisation in Danish municipal data, while \citet{nielsson2024} estimate a precise zero on later Danish variation and attribute it to pass-through to tenants under rent regulation. Our capitalisation figures should therefore be read as the textbook upper bound on the wealth effect, not as a forecast. Section~\ref{sec:robustness} quantifies the first of these: full capitalisation at discount rates of 3--5 per cent implies a fall in land values of 15--26 per cent, or \pounds 1.15--1.92 trillion, concentrated on the top wealth decile. Our flow-based accounting is best interpreted as the impact incidence in the initial years of the reform, and it understates how sharply the reform bears on land wealth.

\paragraph{Regional land shares.} Equation~(\ref{eq:land}) applies a single land share $\lambda_r$ to every property in a region. In practice land shares vary within regions --- they are higher in urban centres and lower for new-build dwellings on inexpensive land --- so we are likely to understate within-region heterogeneity in LVT liabilities, although decile-level averages should be less affected.

\paragraph{Corporate land allocation.} We allocate \pounds 2.06 trillion of corporate land to households in proportion to their measured corporate wealth. This abstracts from foreign ownership of UK corporate equity, which would shift part of the burden abroad and render the reform more favourable to UK households than we report, and it assumes full pass-through of the corporate-sector tax to shareholders rather than to workers or consumers.

\paragraph{Northern Ireland.} Northern Ireland levies domestic rates rather than council tax; the central specification leaves rates unchanged, so Northern Irish households pay LVT in addition to their existing property tax. Section~\ref{sec:robustness} reports the Great Britain-only alternative, which raises the required rate to 0.79 per cent and leaves the distributional pattern intact. A comprehensive UK-wide reform would fold domestic rates into the replacement, lowering the required rate slightly below our central figure.

\paragraph{Measurement.} Land values are imputed rather than observed: Wealth and Assets Survey property wealth is self-reported, the upper tail of the survey remains thin even after enhancement, and the National Balance Sheet's land estimates rest on the ONS's residual valuation method \citep{ons2025}. Council tax liabilities are modelled from regional band averages rather than parcel-level records.

\subsection{Policy implications and design options}

The most salient losing group comprises asset-rich, income-poor households --- typically older owner-occupiers of valuable homes, concentrated in London and the South East --- for whom an annual charge of 0.77 per cent of land value may represent a substantial share of cash income. Established remedies address this liquidity problem: deferral of the liability, with interest, until sale or death, as operated for property taxes in several US states and recommended by \citet{mirrlees2011} and \citet{muellbauer2018}; a phased transition over an extended period; or an exempt band of land value per household. Each option sacrifices some revenue or efficiency, but the underlying difficulty is one of liquidity rather than solvency, since the households concerned are by construction wealthy. The appendix quantifies the first option: households containing someone of state pension age owe 35.1 per cent of the total yield, \pounds 20.3 billion, so a universal pensioner deferral would postpone roughly a third of first-year receipts. That is the scale of the transitional financing requirement --- large, bounded, and a cash-flow cost rather than a revenue loss, since deferred liabilities are recovered with interest from estates.

An alternative reform family would replace council tax with a proportional tax on total property value (structures plus land), as proposed by \citet{muellbauer2005} and modelled by \citet{corlett2018}. A proportional property tax shares the LVT's correction of council tax's regressivity in value and obsolete valuations, and is administratively simpler because market prices are observed; however, it retains the disincentive to improve property and taxes the elastic (structures) margin alongside the inelastic one. We simulate this comparator directly in Section~\ref{sec:robustness}. At the 0.73 per cent rate that raises the same \pounds 57.6 billion, a proportional property tax produces fewer winners (64.4 against 68.2 per cent) and a shallower wealth gradient: the top wealth decile loses \pounds 4{,}217 rather than \pounds 5{,}752. The land base is more concentrated than the property base because location value is more concentrated than structure value, so for a given revenue target the LVT asks more of the largest landholders and less of everyone else. That is the quantitative content of the land/structures distinction, and it is the price paid for the proportional property tax's administrative advantage of observable market prices.

\subsection{Conclusion}

This paper has presented, to our knowledge, the first open-source household-level microsimulation of a revenue-neutral replacement of council tax with a land value tax for the United Kingdom. At the budget-neutral rate of 0.77 per cent, our results suggest that the reform would benefit roughly two-thirds of households, reduce after-housing-cost poverty modestly, and reallocate the burden of recurrent property taxation from council tax's regressive banded structure to the distribution of land wealth, with the largest losses concentrated in the top wealth decile. Measured income inequality and before-housing-cost poverty would be essentially unaffected --- a result that underscores the importance of examining incidence along the wealth as well as the income distribution, and a caution against reading the near-zero movement in headline statistics as evidence that little has changed. The principal obstacles to implementation are practical rather than analytical: valuing land separately from structures at scale, and managing the transition for asset-rich, income-poor owners. Both problems are long-standing, and both admit established policy responses. Extending the analysis to incorporate capitalisation, behavioural responses and parcel-level land valuation remains a subject for future research.


\begin{thebibliography}{99}

\bibitem[Adam et~al.(2020)]{adam2020}
Adam, S., Hodge, L., Phillips, D. and Xu, X. (2020).
\newblock Revaluation and reform: bringing council tax in England into the 21st century.
\newblock IFS Report R168, Institute for Fiscal Studies, London.

\bibitem[ap Gwilym et~al.(2020)]{apgwilym2020}
ap Gwilym, R., Jones, E. and Rogers, H. (2020).
\newblock A technical assessment of the potential for a local land value tax in Wales.
\newblock Bangor University, report for the Welsh Government, Cardiff.

\bibitem[Arnott and Stiglitz(1979)]{arnott1979}
Arnott, R.~J. and Stiglitz, J.~E. (1979).
\newblock Aggregate land rents, expenditure on public goods, and optimal city size.
\newblock \emph{The Quarterly Journal of Economics}, 93(4), 599--616.

\bibitem[Bonnet et~al.(2021)]{bonnet2021}
Bonnet, O., Chapelle, G., Trannoy, A. and Wasmer, E. (2021).
\newblock Land is back, it should be taxed, it can be taxed.
\newblock \emph{European Economic Review}, 134, 103696.

\bibitem[Corlett and Gardiner(2018)]{corlett2018}
Corlett, A. and Gardiner, L. (2018).
\newblock Home affairs: options for reforming property taxation.
\newblock Resolution Foundation, London.

\bibitem[Dye and England(2009)]{dye2009}
Dye, R.~F. and England, R.~W. (eds.) (2009).
\newblock \emph{Land Value Taxation: Theory, Evidence, and Practice}.
\newblock Lincoln Institute of Land Policy, Cambridge, MA.

\bibitem[Feldstein(1977)]{feldstein1977}
Feldstein, M. (1977).
\newblock The surprising incidence of a tax on pure rent: a new answer to an old question.
\newblock \emph{Journal of Political Economy}, 85(2), 349--360.

\bibitem[George(1879)]{george1879}
George, H. (1879).
\newblock \emph{Progress and Poverty: An Inquiry into the Cause of Industrial Depressions and of Increase of Want with Increase of Wealth}.
\newblock D. Appleton and Company, New York.

\bibitem[Murphy and Snelling(2019)]{ippr2019}
Murphy, L. and Snelling, C. (2019).
\newblock A poor tax: reforming council tax in London --- final report.
\newblock Institute for Public Policy Research, London, May 2019.

\bibitem[Fairer Share(2021)]{fairershare2021}
Fairer Share and WPI Economics (2021).
\newblock The proportional property tax: a fairer way to tax property in Britain.
\newblock Fairer Share, London.

\bibitem[Hughes et~al.(2018)]{hughes2018}
Hughes, C., Sayce, S., Shepherd, E. and Wyatt, P. (2018).
\newblock Investigation of potential land value tax policy options for Scotland: final report.
\newblock Scottish Land Commission, Inverness.

\bibitem[Johansson et~al.(2008)]{johansson2008}
Johansson, {\AA}., Heady, C., Arnold, J., Brys, B. and Vartia, L. (2008).
\newblock Taxation and economic growth.
\newblock OECD Economics Department Working Paper No.~620, OECD Publishing, Paris.

\bibitem[Leunig(2024)]{leunig2024}
Leunig, T. (2024).
\newblock A fairer property tax.
\newblock Onward, London.

\bibitem[Mirrlees et~al.(2011)]{mirrlees2011}
Mirrlees, J., Adam, S., Besley, T., Blundell, R., Bond, S., Chote, R., Gammie, M., Johnson, P., Myles, G. and Poterba, J. (2011).
\newblock \emph{Tax by Design: The Mirrlees Review}.
\newblock Oxford University Press, Oxford.

\bibitem[Muellbauer(2005)]{muellbauer2005}
Muellbauer, J. (2005).
\newblock Property taxation and the economy after the Barker Review.
\newblock \emph{The Economic Journal}, 115(502), C99--C117.

\bibitem[Muellbauer(2018)]{muellbauer2018}
Muellbauer, J. (2018).
\newblock Housing, debt and the economy: a tale of two countries.
\newblock \emph{National Institute Economic Review}, 245(1), R20--R33.

\bibitem[Oates and Schwab(1997)]{oates1997}
Oates, W.~E. and Schwab, R.~M. (1997).
\newblock The impact of urban land taxation: the Pittsburgh experience.
\newblock \emph{National Tax Journal}, 50(1), 1--21.

\bibitem[OBR(2025)]{obr2025}
Office for Budget Responsibility (2025).
\newblock Economic and Fiscal Outlook, November 2025.
\newblock OBR, London.

\bibitem[ONS(2025)]{ons2025}
Office for National Statistics (2025).
\newblock National balance sheet estimates for the UK: 2025.
\newblock ONS, Newport. \url{https://www.ons.gov.uk/economy/nationalaccounts/uksectoraccounts/bulletins/nationalbalancesheet/2025}.

\bibitem[PolicyEngine(2026)]{policyengine}
PolicyEngine (2026).
\newblock PolicyEngine UK: an open-source microsimulation model of the UK tax and benefit system.
\newblock \url{https://policyengine.org/uk}.

\bibitem[OECD(2024)]{oecd2024}
OECD (2024).
\newblock OECD Economic Surveys: United Kingdom 2024.
\newblock OECD Publishing, Paris.

\bibitem[Ramsey(1927)]{ramsey1927}
Ramsey, F.~P. (1927).
\newblock A contribution to the theory of taxation.
\newblock \emph{The Economic Journal}, 37(145), 47--61.

\bibitem[Tideman(1995)]{tideman1995}
Tideman, N. (1995).
\newblock Taxing land is better than neutral: land taxes, land speculation and the timing of development.
\newblock Lincoln Institute of Land Policy Working Paper, Cambridge, MA. Reprinted in K.~C. Wenzer (ed.), \emph{Land-Value Taxation}, M.~E. Sharpe, 1999.

\bibitem[Scottish Government(2015)]{scotgov2015}
Commission on Local Tax Reform (2015).
\newblock Just change: a new approach to local taxation.
\newblock Scottish Government, Edinburgh.

\end{thebibliography}

\appendix
\section{Appendix}
\label{sec:appendix}

\subsection{Impact by income decile at alternative rates}

Table~\ref{tab:rategrid} reports the average net income change by income decile across the full rate grid, with council tax abolished in every scenario. At 0.5 per cent every decile except the top gains on average, and the top is essentially unchanged ($-$\pounds 7), the reform reducing net revenue by \pounds 21.3 billion; from 1.5 per cent upwards every decile loses on average, with losses concentrated in deciles 9--10. Rates other than 0.79 per cent are not revenue-neutral, so these columns combine the incidence of the swap with that of a net tax cut or rise.

\begin{table}[htbp]
\centering
\caption{Average net income change (\pounds/year) by income decile and LVT rate}
\label{tab:rategrid}
\begin{tabular}{lrrrrrrr}
\toprule
Decile & 0.5\% & 0.79\% & 1.0\% & 1.5\% & 2.0\% & 3.0\% & 5.0\% \\
\midrule
1  & $+$906     & $+$584     & $+$344     & $-$218     & $-$780     & $-$1{,}903  & $-$4{,}150 \\
2  & $+$568     & $-$17      & $-$454     & $-$1{,}477 & $-$2{,}500 & $-$4{,}545  & $-$8{,}636 \\
3  & $+$929     & $+$640     & $+$425     & $-$80      & $-$584     & $-$1{,}593  & $-$3{,}611 \\
4  & $+$639     & $+$122     & $-$264     & $-$1{,}168 & $-$2{,}071 & $-$3{,}877  & $-$7{,}490 \\
5  & $+$562     & $-$87      & $-$573     & $-$1{,}707 & $-$2{,}841 & $-$5{,}110  & $-$9{,}647 \\
6  & $+$789     & $+$173     & $-$288     & $-$1{,}365 & $-$2{,}442 & $-$4{,}595  & $-$8{,}903 \\
7  & $+$931     & $+$385     & $-$24      & $-$978     & $-$1{,}933 & $-$3{,}842  & $-$7{,}659 \\
8  & $+$894     & $+$161     & $-$387     & $-$1{,}668 & $-$2{,}950 & $-$5{,}512  & $-$10{,}637 \\
9  & $+$387     & $-$668     & $-$1{,}458 & $-$3{,}302 & $-$5{,}146 & $-$8{,}835  & $-$16{,}212 \\
10 & $-$7       & $-$1{,}429 & $-$2{,}493 & $-$4{,}979 & $-$7{,}465 & $-$12{,}437 & $-$22{,}382 \\
\bottomrule
\end{tabular}

\medskip
{\footnotesize\emph{Note:} Rates above 0.79\% raise net revenue that is not recycled; rates below it reduce net revenue. \emph{Source:} Authors' calculations using PolicyEngine UK.}
\end{table}

\subsection{Landless households}

Households holding no land --- no residential property, directly owned land or corporate wealth --- make up 17.2 per cent of households. Because their LVT liability is zero at any rate, their outcome is invariant across the rate grid: an average gain of \pounds 885 per year (their abolished council tax bill), with 77.1 per cent winning. The non-winners are predominantly households whose council tax was already fully covered by council tax reduction, for whom abolition changes nothing.

\subsection{Worked single-household example}

Consider a household owning a \pounds 400{,}000 home in London, where the land share is 85 per cent, implying household land of \pounds 340{,}000. If the household also holds \pounds 50{,}000 of corporate wealth, its attributed corporate land is that sum multiplied by the population ratio of aggregate corporate land to aggregate corporate wealth, which the simulation puts at \pounds 0.14514 per pound of corporate wealth, or \pounds 7{,}257. Its total land value is therefore \pounds 347{,}257 and its annual liability at the budget-neutral rate (0.786 per cent before rounding) is \pounds 2{,}730; at a 1 per cent rate it would be \pounds 3{,}473. Whether the household gains depends on the abolished council tax bill: with an illustrative upper-band London bill of \pounds 2{,}400 it loses roughly \pounds 330 a year, and it breaks even at a bill of \pounds 2{,}730.

The allocation ratio must be computed against the weighted population. Evaluating equation~(\ref{eq:land}) inside a single-household simulation assigns the entire \pounds 2.06 trillion corporate land base to that one record, because it constitutes the whole population; the example above therefore uses the population ratio from the microsimulation.

\subsection{Capitalisation by wealth decile}

Table~\ref{tab:capitalisation} reports the one-off wealth loss implied by full capitalisation of the permanent 0.79 per cent levy, $\Delta W_i = -(r^{*}/\rho) L_i$, at discount rates of 3, 4 and 5 per cent (the table reports the 3 and 5 per cent columns), under the assumption that the levy is assessed for ever on the fixed pre-reform base --- a levy reassessed on post-capitalisation values falls by $r^{*}/(\rho+r^{*})$ instead, so these are upper bounds. The aggregate loss is \pounds 1.17 trillion at $\rho=5$ per cent, \pounds 1.46 trillion at 4 per cent and \pounds 1.95 trillion at 3 per cent, corresponding to falls in land prices of 15.7, 19.7 and 26.2 per cent. Expressed as a share of total wealth the loss peaks in the middle of the wealth distribution --- deciles 5 and 6, whose wealth is overwhelmingly housing --- while in absolute terms it is concentrated at the top.

\begin{table}[htbp]
\centering
\caption{Capitalised one-off wealth loss by wealth decile (\pounds), 0.79\% rate}
\label{tab:capitalisation}
\begin{tabular}{lrrrr}
\toprule
Wealth & \multicolumn{2}{c}{$\rho = 3\%$} & \multicolumn{2}{c}{$\rho = 5\%$} \\
\cmidrule(lr){2-3}\cmidrule(lr){4-5}
decile & Loss (\pounds) & \% of wealth & Loss (\pounds) & \% of wealth \\
\midrule
1  & 0           & ---  & 0           & --- \\
2  & 8           & 2.9  & 5           & 1.7 \\
3  & 329         & 3.7  & 197         & 2.2 \\
4  & 6{,}574     & 8.1  & 3{,}944     & 4.9 \\
5  & 20{,}017    & 9.8  & 12{,}010    & 5.9 \\
6  & 36{,}301    & 9.5  & 21{,}780    & 5.7 \\
7  & 55{,}369    & 8.4  & 33{,}222    & 5.1 \\
8  & 83{,}911    & 8.0  & 50{,}347    & 4.8 \\
9  & 117{,}776   & 6.9  & 70{,}666    & 4.1 \\
10 & 284{,}431   & 6.6  & 170{,}659   & 4.0 \\
\bottomrule
\end{tabular}

\medskip
{\footnotesize\emph{Note:} Loss is $(r^{*}/\rho) L_i$ averaged within decile; ``\% of wealth'' is that loss over average total wealth. Deciles 1--2 hold negligible land.}
\end{table}

\subsection{Deferral for pensioner households}

The liquidity problem the reform creates is concentrated and measurable. Households containing at least one person of state pension age make up 27.2 per cent of households and would owe \pounds 20.6 billion of the \pounds 58.5 billion total, or 35.2 per cent of the yield, at an average liability of \pounds 2{,}354. This is an upper bound on what a deferral scheme could attract: it counts every pensioner household's full liability, including attributed corporate land and households that would not elect to defer, and models no eligibility test, take-up, interest or settlement flows. A full deferral option --- the liability rolled up with interest against the property, as recommended by \citet{mirrlees2011} --- would therefore postpone at most roughly a third of receipts in the first year, declining thereafter as deferred estates are settled. This bounds the transitional financing problem: it is a cash-flow cost to government rather than a revenue loss.

\subsection{Data and parameter summary}

\begin{table}[htbp]
\centering
\footnotesize
\caption{Key data sources and parameters}
\label{tab:params}
\begin{tabular}{@{}llp{6.4cm}@{}}
\toprule
Item & Value & Source / method \\
\midrule
\multicolumn{3}{@{}l}{\emph{Simulation}} \\
\quad Microdata & EFRS 2023--24 & FRS enhanced with SPI quantile-forest imputation and WAS wealth \\
\quad Model version & policyengine-uk 2.89.4 & Released version used; see replication note below \\
\quad Simulation year & 2026--27 & Uprated from the 2023--24 survey \\
\quad Households & 32.08m & Weighted EFRS household count \\
\addlinespace
\multicolumn{3}{@{}l}{\emph{Council tax baseline (replacement target)}} \\
\quad Gross liabilities & \pounds 61.9bn & Modelled, VOA band calibration \\
\quad Council tax reduction & \pounds 3.4bn & Modelled \\
\quad Net revenue ($R_{\mathrm{CT}}$) & \pounds 58.5bn & Equation~(\ref{eq:ctrev}) \\
\quad \emph{Memo:} OBR receipts & \pounds 53.7bn & OBR EFO; robustness row \\
\addlinespace
\multicolumn{3}{@{}l}{\emph{Land base}} \\
\quad Household land & \pounds 5.38tn & Model output, equation~(\ref{eq:land}) \\
\quad Corporate land & \pounds 2.06tn & ONS NBS aggregate, allocated by corporate wealth \\
\quad Total base ($L$) & \pounds 7.44tn & 104.8\% of un-uprated ONS 2024 (\pounds 7.10tn) \\
\quad Regional land shares ($\lambda_r$) & 0.42--0.85 & MHCLG residual method, equation~(\ref{eq:lambda}) \\
\addlinespace
\multicolumn{3}{@{}l}{\emph{Reform}} \\
\quad Budget-neutral rate ($r^{*}$) & 0.79\% & $R_{\mathrm{CT}}/L$ \\
\quad LVT parameter & \multicolumn{2}{l}{\texttt{gov.contrib.ubi\_center.land\_value\_tax.rate}} \\
\quad Abolition parameter & \multicolumn{2}{l}{\texttt{gov.contrib.abolish\_council\_tax}} \\
\bottomrule
\end{tabular}
\end{table}

\subsection{Replication note}

All results are generated from a single baseline simulation of PolicyEngine UK 2.89.4 on the pinned Enhanced FRS release \texttt{enhanced\_frs\_2023\_24@1.56.14}, together with the closed-form arithmetic validated in Section~\ref{sec:robustness}. The generating script (\texttt{uk\_lvt/pipeline\_direct.py}) re-runs the full model at three rates on every invocation and refuses to write results if the arithmetic and the model disagree by more than \pounds 1 per household or 0.01 percentage points of poverty.

Three caveats on measurement and version sensitivity are worth recording. First, all poverty rates are shares of individuals (household poverty status weighted by household size, the HBAI convention) and the income Gini is person-weighted over equivalised HBAI net income. Second, poverty \emph{levels} move with the underlying benefit and uprating parameters across model and data releases; under the pinned releases the budget-neutral changes are $-0.23$ points BHC and $-0.69$ points AHC. Third, the wealth-decile results depend on how deciles are constructed in the presence of tied zero-wealth households; we construct equal-weight deciles directly rather than using the model's own decile variable.


\end{document}
